\section{Conclusion}

This paper established a closure cost model for the 15 highest-risk interstate corridors
in the ROUTE corpus and derived redundancy value as the economic metric for alternate
route investment justification. Three findings are central.

\textbf{The top-5 closure-cost corridors account for \$5.4 billion in annual freight
disruption, but none are the highest V/C corridors.} Donner Pass (\$2.4B) is the most
expensive closed corridor in the national network despite operating well below capacity
when open. Gulf Coast I-10 (\$1.2B) is expensive because of the high freight volume
(12,000 trucks/day, amplified by New Orleans port activity) combined with the long
mean closure duration (48 hours) for hurricane events. Dallas (\$0.8B) is the only
urban interchange in the top-5, where high truck volume partially compensates for lower
isolation penalty. The V/C metric does not identify any of these corridors as high
priority; the closure cost model does.

\textbf{Donner Pass costs 4.2$\times$ more per closure event than the Dallas
Interchange}, despite lower peak V/C. The isolation multiplier (B1 penalty: 550-mile
I-40 detour vs.\ 200-mile I-20 detour for Dallas) and the higher closure frequency
(50/yr vs.\ 6/yr) combine to produce this differential. The per-closure-hour cost
is actually higher at Dallas (high truck volume), but Donner's 900 annual closure-hours
vastly exceed Dallas's 48. The interaction between isolation, frequency, and volume
is multiplicative --- it cannot be captured by any single-dimension metric.

\textbf{Redundancy value quantifies what alternate routes are worth.} The Donner
redundancy value (\$1.9B/yr from a hypothetical I-70W alternate) is large enough to
justify major capital investment at standard federal discount rates --- but the snowshed
program (\$0.8B capital, \$0.96B/yr benefit, 0.83-year payback, NPV \$+11.1B) is
the correct first investment because it addresses the highest-frequency events at a
fraction of the capital cost.

\textbf{Limitations.} The FHWA incident database undercounts rural closures by an
estimated 30--40\% relative to urban incidents. Donner's \$2.4B annual cost estimate
is therefore conservative: actual costs, including unrecorded short-duration events,
may be 30\% higher. Additionally, the model uses ATRI's \$225/hr truck-hour figure
uniformly; for agricultural commodity corridors (Snoqualmie, I-35 Oklahoma), the
per-hour cost during harvest season is materially higher due to perishable load
penalties, meaning Snoqualmie's \$0.31B/yr estimate is understated during
September--November.

\textbf{Next steps.} The closure cost model developed here ($E[\text{Cost}] =
\lambda V \cdot E[\min(C_w, C_r)]$) is the input to the managed lane NPV model
in E.1 and the Interstate 2.0 framework investment prioritization in E.2
\citep{ROUTE_D1}. Papers E.1 and E.2 use the D.2 closure cost as one component
of a multi-benefit NPV calculation that also includes congestion relief, multimodal
integration, and resilience co-benefits. The D.2 framework establishes that resilience
investment and congestion investment are not competing for the same dollar: they target
different failure modes on different corridors, and the highest-NPV investments in the
ROUTE framework address both simultaneously through compound strategies that neither
PROTECT nor NHPP currently targets.

\subsection*{Policy Recommendation}

FHWA should compute redundancy value for all T1 corridors as a standard investment
criterion alongside V/C ratio and ATRI bottleneck cost. The computation requires:
(1) FHWA incident database closure frequency and duration data, which already exists;
(2) HPMS truck AADT, which already exists; (3) network-based detour distance
computation, which requires a one-time graph analysis. The redundancy value criterion
would redirect a portion of NHPP and PROTECT funding toward the highest-consequence
resilience gaps in the network --- and away from the incremental congestion relief
investments that currently dominate both programs despite substantially lower NPV
per dollar.
